% ============================================================
% Exact nonequilibrium steady state of collective group-flip dynamics
% on quenched interaction neighborhoods
%
% Target: Physical Review E (Regular Article). REVTeX 4.2, apsrev4-2.
% Overleaf: set this file as the main document; compiler pdfLaTeX.
% Files needed in the project: main.tex, refs_pre.bib,
%   fig1_schematic.png, fig3_containment.pdf (or .png),
%   fig2_quenched_annealed.pdf.
% The Supplemental Material is SM.tex (compile separately).
% ============================================================
\documentclass[aps,pre,reprint,superscriptaddress,amsmath,amssymb,longbibliography]{revtex4-2}

\usepackage{graphicx}
\usepackage{dcolumn}
\usepackage{bm}
\usepackage{amsthm}
\usepackage{booktabs}
\usepackage[hyphens]{url}
\usepackage[colorlinks=true,linkcolor=blue,citecolor=blue,urlcolor=blue]{hyperref}

\newtheorem{theorem}{Theorem}
\newtheorem{lemma}{Lemma}
\newtheorem{proposition}{Proposition}
\newtheorem{corollary}{Corollary}
\newtheorem{conjecture}{Conjecture}
\theoremstyle{remark}
\newtheorem{remark}{Remark}

\newcommand{\ctil}{\tilde c}
\newcommand{\cG}{c_G}
\newcommand{\E}{\mathbb{E}}

\begin{document}

\title{Exact nonequilibrium steady state of collective group-flip dynamics\\
on quenched interaction neighborhoods}

\author{Pham Le Dat Nguyen}
\email{phmldtngyn.777@gmail.com}
\affiliation{Independent Researcher, Ha Noi, Viet Nam}

\date{\today}

\begin{abstract}
Group-update models of opinion and spin dynamics are usually analyzed on
annealed interaction structures, where the interacting group is redrawn
at every step, and they relax to absorbing consensus. We study a
two-state group-flip dynamics in which a firing node and every node that
shares its state inside its fixed (quenched) graph-metric neighborhood
switch state simultaneously. The dynamics has no absorbing state and no
external noise, so its long-time behavior is a genuine nonequilibrium
steady state. If every neighborhood is a strict majority of the system
(R1), every complement of a neighborhood lies inside the neighborhood of
each of its own nodes (R2), and all neighborhoods have $\beta$ nodes, we
prove that the stationary majority fraction is
$\E[M]=(2-q)/(3-2q)$ with $q=\beta/n$, exactly for every system size $n$
and independently of dimension and topology; equivalently, the mean
absolute magnetization is $1/(3-2q)$. The proof rests on an exact affine
closure for the mean minority size. Under R1 alone we prove that
$\E[M]$ equals the fraction of updates that reverse the sign of the
magnetization, and that the steady state violates detailed balance
through one-way minority-erasure transitions, whose rate is
$n(1-q)/(3-2q)$ under R1 and R2. For unequal neighborhoods the closure
survives with one effective parameter and gives exact two-sided bounds.
An annealed version with the same neighborhood size orders less in every
case computed with $n-\beta\ge2$ (up to $n=3200$), so the exact steady
state is a property of quenched structure. Below the R2 threshold, exact solutions of all 17
ring cases with $n\le33$ lie below the closed form, which we conjecture
to be an upper bound there.
\end{abstract}

\maketitle

% ============================================================
\section{Introduction}
\label{sec:intro}
% ============================================================

Stochastic models of interacting two-state units, such as the voter
model and majority-rule dynamics, are standard testing grounds of
nonequilibrium statistical
physics~\cite{CliffordSudbury1973,HolleyLiggett1975,deOliveira1992,Krapivsky2003,Castellano2009,Redner2019}.
On a finite connected graph the voter model reaches one of its two
absorbing consensus states, and a nontrivial stationary state requires
an additional mechanism, typically spontaneous state changes (noise), as
in the noisy voter or Kirman
model~\cite{Kirman1993,GranovskyMadras1995,Carro2016}. Updates that
involve groups of units, such as majority rule and nonlinear $q$-voter
rules, introduce nonlinear drift terms that are absent from dyadic
copying~\cite{Krapivsky2003,Castellano2009qvoter,Neuhauser2020,Sahasrabuddhe2021,Noonan2021},
and group interactions are now studied systematically on higher-order
structures~\cite{Battiston2020,Iacopini2019}.

Many analytical results for group updates rely on annealed or
mean-field structures, in which the interacting group is redrawn at every step, so
that a mean-field drift can be written. A recent example is the
group-driven voter model (GVM) of Kim \emph{et al.}~\cite{Kim2025}: on
annealed uniform hypergraphs the time to consensus grows as
$A\ln N$, and for hyperedges larger than three it is shortest at an
intermediate nonlinearity. Annealing
is also a common route to analytical
treatments of pairwise voter dynamics on complex
networks~\cite{Vilone2004,Carro2016}. Two questions are then natural.
Can a group-update dynamics on a \emph{quenched} structure be solved
exactly? And if it can, does the quenched structure change the
macroscopic state, or does an annealed description with the same group
size reproduce it?

Here we answer both questions for a simple rule that we call group-flip
dynamics. Each node fires at rate one; when node $i$ fires, the node and
every node that shares its state inside a fixed metric ball $B_i(r)$
switch state together. Two features distinguish it from the models
above. First, the update moves an entire same-state cluster of a
neighborhood at once and involves no copying of a neighbor's state and
no unanimity condition. Second, no state is absorbing: from consensus,
the next fire creates a minority. The long-time state is therefore a
stationary state without external noise, and we show that it violates
detailed balance. The cluster move is reminiscent of cluster Monte Carlo
algorithms~\cite{SwendsenWang1987}, with the difference that there
clusters are built to preserve a Boltzmann measure, whereas here the
move is unconditional and the stationary measure is not reversible.

Our results are as follows.
(i)~Under a strict-majority condition R1 on the balls, the sign of the
magnetization reverses exactly when a majority node fires; the
stationary majority fraction $\E[M]$ therefore equals the fraction of
updates that reverse the sign, and the dynamics contains one-way
transitions, so the steady state is a nonequilibrium steady state (NESS)
(Sec.~\ref{sec:R1}).
(ii)~If, in addition, every residual (the complement of a ball) lies
inside the ball of each of its own nodes (R2) and all balls have the
same size $\beta$, then
\begin{equation}
\E[M]=\frac{2-q}{3-2q},\qquad q=\frac{\beta}{n},
\label{eq:main}
\end{equation}
exactly for every finite $n$ and for every geometry satisfying the
conditions (Sec.~\ref{sec:band}). The derivation rests on an exact
affine closure of the conditional mean minority size, so that no
moment-closure approximation is involved. Related observables (mean
absolute magnetization, rotation rate, erasure current, cycle-length
law) follow in closed form, and for unequal balls the closure survives
with a single effective parameter and yields exact two-sided bounds.
(iii)~An annealed version with the same ball size, solved exactly as a
Markov chain on the magnetization, orders less than the quenched system
in every case computed with $n-\beta\ge2$, up to $n=3200$; the gap does not close as $n$
grows in that range. Equation~\eqref{eq:main} is thus a property of
quenched neighborhoods, and the exact lemmas identify the mechanism
(Sec.~\ref{sec:annealed}).
(iv)~Below the R2 threshold, minorities are no longer erased in one
update. We give the one-step extinction probability in closed form on
rings, solve the stationary chain for all 17 ring cases with $n\le33$,
and find that all values lie below Eq.~\eqref{eq:main}, which we state
as a conjecture (Sec.~\ref{sec:regionII}).
Table~\ref{tab:compare} places the model among related dynamics.

\begin{table*}[t]
\caption{Group-flip dynamics compared with related two-state dynamics.
The entries describe the settings analyzed in the cited works; ``long-time
state'' refers to a finite system.}
\label{tab:compare}
\begin{ruledtabular}
\begin{tabular}{p{3.2cm}p{4.1cm}p{4.1cm}p{4.1cm}}
Model & Interaction structure analyzed & Update & Long-time state\\
\colrule
Voter model~\cite{CliffordSudbury1973,HolleyLiggett1975} & graph & node copies a random neighbor & absorbing consensus\\
Noisy voter (Kirman) model~\cite{Kirman1993,GranovskyMadras1995,Carro2016} & complete graph; networks (annealed approx.) & copying plus spontaneous flips & stationary, sustained by noise\\
Majority rule~\cite{Krapivsky2003,Noonan2021} & mean field; hypergraphs & group adopts its local majority & absorbing consensus\\
Group-driven voter model~\cite{Kim2025} & annealed uniform hypergraphs & node adopts a unanimous sample & absorbing consensus\\
This work & quenched metric balls & same-state part of a ball switches & stationary without noise; exact $\E[M]$\\
\end{tabular}
\end{ruledtabular}
\end{table*}

% ============================================================
\section{Model}
\label{sec:model}
% ============================================================

\subsection{Quenched neighborhoods}

Let $G=(V,E)$ be a connected graph with $n$ nodes and graph distance
$d$. For a radius $r$, the interaction neighborhood of node $i$ is the
ball $B_i=\{j:d(i,j)\le r\}$, of size $\beta_i=|B_i|$; when all balls
have the same size we write $\beta$ and define the coverage
$q=\beta/n$. The balls are quenched: they are fixed for the whole
dynamics. Since $d$ is symmetric, $j\in B_i\iff i\in B_j$, a property
used repeatedly below. The complement $R_i=V\setminus B_i$ is the
residual of $i$.

Our main example is the ring $C_n(1,2)$, in which each node is linked to
its first and second neighbors on a cycle. Its distance is
$d(i,j)=\lceil\delta(i,j)/2\rceil$, with $\delta$ the circular lattice
distance, so the balls are arcs of $\beta=\min(n,4r+1)$ nodes. In two
dimensions we use the analogous square lattice with first- and
second-neighbor links along each axis, for which
$\beta=8r^2+1$ without wrap-around~\cite{SM}.

\begin{remark}[Origin of the construction]
\label{rem:2section}
These lattices are the 2-section graphs of uniform hypergrids whose
hyperedges are three consecutive nodes along an axis, the formulation
used in an earlier version of this work~\cite{repo}. The dynamics
depends on a hypergraph only through its 2-section graph, so no
higher-order interaction beyond the pairwise projection is involved,
and we use graph language throughout.
\end{remark}

\subsection{Group-flip dynamics}

Each node carries a state $\sigma_i\in\{0,1\}$ and fires at rate one,
independently. When node $i$ fires, every node $j\in B_i$ with
$\sigma_j=\sigma_i$ switches to $1-\sigma_i$ simultaneously
(Fig.~\ref{fig:schematic}). A single update therefore moves the
same-state part of a ball, and its outcome is deterministic given the
firing node. The total update rate is $n$ in every configuration. The
version with $Q>2$ states, in which the moved class adopts a state
chosen uniformly among the other $Q-1$, is solved exactly in the
well-mixed limit $B_i=V$ in the Supplemental Material~\cite{SM}; all
results below are for $Q=2$.

With $n_\sigma$ the number of nodes in state $\sigma$, the
magnetization is $m=(n_1-n_0)/n$ and the majority fraction is
$M=\max(n_0,n_1)/n=(1+|m|)/2$. The primary observable is the stationary
(long-run time-averaged) value $\E[M]$. A rotation is an update that
changes the majority label, i.e., reverses the sign of $m$. We denote by
$K$ the number of updates between consecutive rotations and by $T$ the
corresponding time; a cycle is the interval between consecutive
rotations.

\begin{figure*}[t]
\centering
\includegraphics[width=\textwidth]{fig1_schematic}
\caption{Group-flip update on a ring. (a)~Ball $B_i(r)$ (shaded) and
residual $R_i$. (b)~A fire at $i$: all nodes of $i$'s state inside $B_i$
(red) switch together. (c)~The two outcomes from a strict-majority
state: a majority fire reverses the majority label
(Theorem~\ref{thm:rotation}); a minority fire from a self-contained
minority erases it (Lemma~\ref{lem:epoch}).}
\label{fig:schematic}
\end{figure*}

\subsection{Geometric conditions}

Two conditions organize the analysis:
\begin{description}
\item[R1:] $\beta_i>n/2$ for all $i$: every ball is a strict majority of
the system;
\item[R2:] $\mathrm{diam}(R_i)\le r$ for all $i$: every residual lies
inside the ball of each of its own nodes.
\end{description}
R2 compares two lengths: the diameter of the largest minority that a
single update from consensus can create (a residual) and the interaction
range $r$. On the ring~\cite{SM},
\begin{equation}
\mathrm{R1}\iff2\beta>n,\qquad
\mathrm{R2}\iff\beta=n\ \text{or}\ 3\beta\ge2n-1,
\label{eq:ringconditions}
\end{equation}
so that R1$\wedge$R2 holds asymptotically for $q\ge2/3$ (Region~I), R1
alone for $1/2<q<2/3$ (Region~II), and neither for $q\le1/2$
(Region~III); we use the exact finite-$n$ forms throughout.

Because R1 requires $\beta_i>n/2$, the radius grows with $n$ at fixed
$q$, as $r=\Theta(n^{1/d})$ on $d$-dimensional
lattices~\cite{SM}. Each update thus acts on a finite fraction of the
system, as in mean-field treatments, but on a set that is fixed in
space. The comparison with an annealed version at the same $\beta$ in
Sec.~\ref{sec:annealed} isolates the effect of this structure from the
effect of the large neighborhood size.

% ============================================================
\section{Exact results under R1: sign reversals and broken detailed balance}
\label{sec:R1}
% ============================================================

\begin{theorem}[Rotation identity]
\label{thm:rotation}
Assume R1 (equal balls are not required). From any state with a strict
majority $\mathcal M$ and minority $\mathcal A$, an update is a rotation
if and only if the firing node belongs to $\mathcal M$; the
instantaneous rotation rate is $|\mathcal M|=nM$. Tie states are left at
the first update and never re-entered, and carry zero stationary mass.
\end{theorem}

\begin{proof}
Let the firing node be $i\in\mathcal M$. After the update the class of
the new state is $\mathcal A\cup(\mathcal M\cap B_i)$ and that of the old
state is $\mathcal M\setminus B_i$. Their size difference is
\begin{equation}
|\mathcal A|+|\mathcal M\cap B_i|-|\mathcal M\setminus B_i|
=(2\beta_i-n)+2|\mathcal A\setminus B_i|>0,
\label{eq:margin}
\end{equation}
so the majority label changes. If $i\in\mathcal A$, the minority can only
shrink, and no rotation occurs. Each majority node fires at rate one.
From a tie, Eq.~\eqref{eq:margin} with $|\mathcal M|=|\mathcal A|$ shows
that every update produces a strict majority, and strict-majority states
are mapped to strict-majority states.
\end{proof}

\begin{corollary}[Static--dynamic identity]
\label{cor:duality}
Under R1, the long-run fraction of updates that reverse the sign of $m$
equals $\E[M]$, the long-run rotation rate per unit time is
$n\E[M]$, and $\E[K]=1/\E[M]$.
\end{corollary}

\begin{proof}
Firing epochs form a Poisson process of rate $n$, and the firing node is
uniform and independent of the configuration. By the PASTA property
applied to the firing epochs, the configuration seen by a fire has the
stationary law, so a fire is a rotation with probability $\E[M]$ by
Theorem~\ref{thm:rotation}. If the chain has several closed classes, the
statement holds class by class~\cite{SM}.
\end{proof}

Corollary~\ref{cor:duality} gives $\E[M]$ a second, dynamical meaning:
it is the probability that a randomly chosen update reverses the sign of
the order parameter. Since $\E[M]>1/2$, the sign of $m$ reverses at a
rate of order $n$ per unit time at every system size, while $|m|$ stays
of order one. The order is therefore an order in $|m|$ without a
persistent sign, in contrast with ordered phases in which reversals of
the magnetization become rarer as the system grows.

\begin{proposition}[Broken detailed balance]
\label{prop:irrev}
Assume R1. Let $x$ be a configuration whose minority $\mathcal A$ is
nonempty and let a fire of a node $j\in\mathcal A$ with
$\mathcal A\subseteq B_j$ map $x$ to a consensus configuration $y$. Then
no single update maps $y$ to $x$. Consequently, whenever such erasure
transitions occur with positive stationary rate, the stationary
dynamics is not reversible: the steady state carries nonzero probability
currents~\cite{Schnakenberg1976}.
\end{proposition}

\begin{proof}
From a consensus configuration, a fire at $k$ switches exactly $B_k$, so
the configurations reachable in one step are those in which the
nonconsensus state occupies a ball. In $x$ that state occupies
$\mathcal A$, and $|\mathcal A|<n/2<\beta_k$ for every $k$ by R1. A
reversible Markov chain requires $w(x\to y)>0\iff w(y\to x)>0$ for all
pairs; this fails.
\end{proof}

By the same PASTA argument, the stationary rate of minority fires is
$n(1-\E[M])$. Where every minority fire is an erasure, as in Region~I
(Lemma~\ref{lem:epoch} below), this is the rate of one-way transitions;
Eq.~\eqref{eq:observables} gives it in closed form.

% ============================================================
\section{Exact steady state under R1 and R2}
\label{sec:band}
% ============================================================

\subsection{Epoch chain and self-containment}

Let $A_m$ be the minority set immediately after the $m$th rotation,
$a_m=|A_m|$, and $i_m$ the node whose fire produced that rotation.

\begin{lemma}[Epoch invariant]
\label{lem:epoch}
Under R1 and R2, $A_m\subseteq R_{i_m}$ for every $m$. Consequently
$A_m\subseteq B_j$ for every $j\in A_m$ (self-containment), so any
minority fire erases the whole minority, and $a_m\le n-\beta_{i_m}$.
Equal balls are not required.
\end{lemma}

\begin{proof}
From consensus a fire at $i$ switches $B_i$, which is the new majority
by R1, so the new minority is $R_i$. Suppose $A\subseteq R_i$ is the
current minority. R2 gives $R_i\subseteq B_j$ for every $j\in R_i$, hence
$A\subseteq B_j$ for every $j\in A$. If the next fire is at $j\in A$, the
whole minority switches, the system reaches consensus, and the following
fire, at some $i'$, creates the minority $R_{i'}$ and a rotation. If the
next fire is at $i\in\mathcal M$, only majority nodes inside $B_i$ switch
(minority nodes in $B_i$ already carry the target state); a rotation
occurs (Theorem~\ref{thm:rotation}) and the new minority is
$\mathcal M\setminus B_i\subseteq R_i$. Both branches preserve the
invariant.
\end{proof}

Lemma~\ref{lem:epoch} is where the quenched geometry enters: every
minority is a compact domain contained in a residual, and it is erased
in one update by any of its own nodes.

\subsection{Exact affine closure}

\begin{lemma}[Overlap identity]
\label{lem:overlap}
Under R1, R2 and equal balls, for $i$ uniform in the majority $\mathcal
M$ and a minority $A$ of size $a$,
\begin{equation}
\E\big[|A\cap B_i|\,\big|\,a\big]=\frac{a(\beta-a)}{n-a}.
\label{eq:overlap}
\end{equation}
\end{lemma}

\begin{proof}
Exchange the order of summation and use $j\in B_i\iff i\in B_j$:
$\sum_{i\in\mathcal M}|A\cap B_i|=\sum_{j\in A}|B_j\cap\mathcal M|$. By
self-containment $B_j\supseteq A$, so $|B_j\cap\mathcal M|=\beta-a$.
Divide $a(\beta-a)$ by $|\mathcal M|=n-a$.
\end{proof}

A minority fire has probability $a/n$; it erases the minority, and the
following fire creates a minority of size $n-\beta$. A majority fire has
probability $1-a/n$ and creates the minority $\mathcal M\setminus B_i$,
of size $n-\beta-a+|A\cap B_i|$. By Lemma~\ref{lem:overlap},
\begin{equation}
\E[a'\mid a,\ \text{majority fire}]=\frac{(n-\beta)(n-2a)}{n-a},
\label{eq:branch}
\end{equation}
and weighting the two branches gives
\begin{equation}
\E[a_{m+1}\mid a_m=a]=(1-q)(n-a).
\label{eq:affine}
\end{equation}
The conditional mean is affine in $a_m$, so taking expectations against
any distribution of $a_m$ produces no correction from higher moments.
This is the step on which the exact result rests.

\subsection{Stationary state}

The epoch sets form a Markov chain on a finite family of subsets. Let
$\pi$ be any invariant distribution and $\bar x=\E_\pi[a]/n$; uniqueness
is not needed. Taking $\pi$-expectations in Eq.~\eqref{eq:affine} gives
$\bar x=(1-q)(1-\bar x)$, hence
\begin{equation}
\bar x=\frac{1-q}{2-q}
\label{eq:xbar}
\end{equation}
for every invariant distribution. A cycle contains a second update if
and only if its first update is a minority fire, so $K\in\{1,2\}$ and
$P(K=2)=\bar x$.

\begin{theorem}[Exact stationary majority fraction]
\label{thm:band}
Assume $Q=2$, R1, R2, and equal balls $\beta$. Then, for every finite
$n$ and every initial condition,
\begin{equation}
\E[M]=\frac{2-q}{3-2q},\qquad q=\frac{\beta}{n}.
\label{eq:band}
\end{equation}
\end{theorem}

\begin{proof}
Inter-update times are exponential with mean $1/n$ and independent of the
sequence of firing nodes. During the first interval of a cycle the
majority fraction is $1-a_m/n$; if a second interval occurs
(probability $a_m/n$), the system is in consensus and $M=1$. The
expected time integral of $M$ over a cycle is therefore
$n^{-1}(1-\bar x)+n^{-1}\bar x=1/n$, and the expected cycle length is
$(1+\bar x)/n$. Cycles are stationary but not independent; the ergodic
theorem applied to the stationary epoch chain gives the ratio of these
means on each closed class, and Eq.~\eqref{eq:xbar} makes it the same on
every class (Appendix~\ref{app:ergodic}). Hence
$\E[M]=1/(1+\bar x)=(2-q)/(3-2q)$.
\end{proof}

Theorem~\ref{thm:band} applies on the R1$\wedge$R2 domain, where
$\E[M]$ ranges from $(4n+1)/(5n+2)$ at the ring boundary $3\beta=2n-1$
(tending to $4/5$) to $1$ at $q=1$. Values of the closed form at smaller
$q$ are used below only as a comparator. The result does not depend on
the dimension or topology of $G$: two geometries with the same $q$ that
satisfy R1 and R2 have the same stationary majority fraction. Native-ball
simulations of rings ($n=400,401$) and two-dimensional lattices
($L=12,30$), for which R1 and R2 were verified by breadth-first search,
reproduce Eq.~\eqref{eq:band} within $2.1\times10^{-5}$ (largest
$|z|=0.8$)~\cite{SM}.

\begin{corollary}[Closed-form steady-state observables]
\label{cor:observables}
Under the assumptions of Theorem~\ref{thm:band},
\begin{equation}
\begin{aligned}
\E|m|&=\frac{1}{3-2q}, & \E[K]&=\frac{3-2q}{2-q},\\
J_{\rm rot}&=n\,\frac{2-q}{3-2q}, & P(K=2)&=\frac{1-q}{2-q},\\
J_{\rm erase}&=n\,\frac{1-q}{3-2q}, & \mathrm{CV}_T^2&=\frac{7-8q+2q^2}{(3-2q)^2},
\end{aligned}
\label{eq:observables}
\end{equation}
where $J_{\rm rot}$ is the long-run rotation rate per unit time,
$J_{\rm erase}$ the long-run rate of one-way erasure transitions
(Proposition~\ref{prop:irrev}), and $\mathrm{CV}_T^2$ the squared
coefficient of variation of the cycle time.
\end{corollary}

\begin{proof}
$|m|=2M-1$. $\E[K]$ and $P(K=2)$ follow from Eq.~\eqref{eq:xbar};
$J_{\rm rot}=n\E[M]$ by Corollary~\ref{cor:duality}; every minority
fire is an erasure by Lemma~\ref{lem:epoch}, so
$J_{\rm erase}=n(1-\E[M])$. For $T$, write $T=G_1+BG_2$ with
$G_{1,2}$ exponential of mean $1/n$ and $B$ Bernoulli of mean $\bar x$,
independent of $G_1,G_2$~\cite{SM}.
\end{proof}

The cycle statistics $\E[K]$ and $\mathrm{CV}_T^2$ test the cycle
decomposition independently of $\E[M]$; in simulations of the same
Region-I geometries they agree with Eq.~\eqref{eq:observables} with
largest $|z|=1.9$~\cite{SM}. The erasure current $J_{\rm erase}$ is
strictly positive for $q<1$, so by Proposition~\ref{prop:irrev} the
steady state of Theorem~\ref{thm:band} is a NESS for every $q<1$.

\subsection{Unequal balls}

Rotation and self-containment do not use equal balls; only
Lemma~\ref{lem:overlap} does. For unequal balls the same exchange of
summation gives $\sum_{i\in\mathcal M}|A\cap B_i|=S_A-a^2$ with
$S_A=\sum_{j\in A}\beta_j$, and the conditional mean becomes
\begin{equation}
\E[a'\mid A]=n-\bar\beta-a\Big(1+\frac{\bar\beta}{n}\Big)+\frac{2S_A}{n},
\label{eq:affine-het}
\end{equation}
with $\bar\beta=n^{-1}\sum_i\beta_i$. This is affine in $a$ and in the
single additional statistic $S_A$, so the closure survives with one
effective parameter.

\begin{proposition}[Unequal balls]
\label{prop:het}
Assume $Q=2$, R1 and R2, $\beta_i<n$ for some $i$, and let $\bar q=\bar\beta/n$,
$q_{\min}=\min_i\beta_i/n$, $q_{\max}=\max_i\beta_i/n$. On each closed
class $c$ of the epoch chain, with stationary law $\pi_c$, define
$q^\ast_c=\E_{\pi_c}[S_A]/(n\E_{\pi_c}[a])\in[q_{\min},q_{\max}]$. Then
\begin{equation}
\E_c[M]=g(q^\ast_c),\qquad g(s)=\frac{2+\bar q-2s}{3-2s},
\label{eq:het}
\end{equation}
and, since $g$ is decreasing,
\begin{equation}
g(q_{\max})\le\E[M]\le g(q_{\min}).
\label{eq:het-bounds}
\end{equation}
For equal balls, $q^\ast=\bar q=q$ and Eq.~\eqref{eq:het} reduces to
Eq.~\eqref{eq:band}.
\end{proposition}

\begin{proof}
Averaging Eq.~\eqref{eq:affine-het} over $\pi_c$ gives
$\bar x_c=(1-\bar q)/(2+\bar q-2q^\ast_c)$. The cycle structure
($K\in\{1,2\}$, $P(K=2\mid a)=a/n$) and the renewal--reward argument of
Theorem~\ref{thm:band} are unchanged, so $\E_c[M]=1/(1+\bar x_c)$.
Finally $g'(s)=-2(1-\bar q)/(3-2s)^2\le0$. Full proof in~\cite{SM}.
\end{proof}

We checked Eq.~\eqref{eq:het} on 60 distinct ball systems with unequal
ball sizes satisfying R1 and R2 (rings with $9\le n\le14$ and up to five
random chords, $r\in\{2,3,4\}$), computing $q^\ast$ directly from the
rotation-weighted epoch law of the exact stationary chain rather than
inferring it from $\E[M]$: the identity holds to $1.2\times10^{-15}$, and
the bounds~\eqref{eq:het-bounds} hold in every case~\cite{SM}. In all 60
cases $q^\ast<\bar q$, i.e., minorities sit preferentially on nodes with
smaller balls, and therefore $\E[M]>(2-\bar q)/(3-2\bar q)$, by between
$0.0008$ and $0.013$. We do not have a proof that $q^\ast\le\bar q$ in
general.

% ============================================================
\section{Quenched versus annealed neighborhoods}
\label{sec:annealed}
% ============================================================

To separate the effect of quenched structure from that of neighborhood
size, consider an annealed version: at every fire of $i$ the interaction
set is $\{i\}$ together with $\beta-1$ nodes drawn uniformly without
replacement from $V\setminus\{i\}$. The number of switched nodes is then
$1+h$ with $h$ hypergeometric, and the dynamics reduces to a Markov chain
on $n_1\in\{0,\dots,n\}$, whose stationary law we compute by a direct
linear solve~\cite{SM}.

Figure~\ref{fig:annealed} compares the two. In Region~II the annealed
value lies below the exact quenched ring value in all 17 cases of
Table~\ref{tab:exact}, by between $0.021$ and $0.026$. In Region~I, at
$n=400$, the annealed value lies below Eq.~\eqref{eq:band} for every
$\beta\le n-2$, by up to $0.022$; at $\beta=n-1$ the two coincide,
since each interaction set then misses a single node. The gap does not
close with system size in the range computed: at $q=0.7$ it is $0.0173$,
$0.0175$, $0.0176$, $0.0177$, $0.0177$ and $0.0178$ for
$n=100,200,400,800,1600$ and $3200$, and at $q=0.8$ it grows from
$0.0068$ ($n=100$) to $0.0071$ ($n=1600$)~\cite{SM}. These are exact
finite-$n$ values; we do not prove a large-$n$ limit.

\begin{figure}[t]
\centering
\includegraphics[width=\columnwidth]{fig2_quenched_annealed}
\caption{Quenched versus annealed neighborhoods of equal size. Solid
black: exact result Eq.~\eqref{eq:band}, valid for any geometry
satisfying R1 and R2 (shaded, asymptotic location $q\ge2/3$); dotted:
the same function outside its domain, used as a comparator. Circles:
exact stationary values on the quenched ring in Region~II
(Table~\ref{tab:exact}). Squares: annealed version with the same
$(n,\beta)$. Red line: annealed version at $n=400$, every integer
$\beta$, exact chain.}
\label{fig:annealed}
\end{figure}

The exact lemmas identify where the difference comes from. Write
$x=a/n$ for the minority fraction. In the quenched system under R1 and
R2, a minority fire erases the minority ($x\to0$), and a majority fire
leaves, on average, $x'=(1-q)(1-2x)/(1-x)$ by Eq.~\eqref{eq:branch}. In
the annealed system at large $n$, hypergeometric concentration gives
$x\to(1-q)x$ for a minority fire and $x\to(1-q)(1-x)$ for a majority
fire. Both quenched branches leave a smaller minority, since
$(1-2x)/(1-x)<1-x$ for $0<x<1/2$. The origin is self-containment: a
minority is a compact domain whose nodes lie in one another's balls, so
it is erased completely by a single fire from inside, and it is
under-represented in the balls of majority nodes
[Eq.~\eqref{eq:overlap}], which therefore leave fewer minority nodes
behind. An annealed interaction set carries no such correlation with the
minority. Iterating the annealed large-$n$ map reproduces the exact
annealed chain at the largest sizes computed (at $q=0.7$:
$0.79472\pm0.00002$ from the map, $0.79475$ for the chain at $n=3200$)~\cite{SM}; this map is a heuristic for
the annealed system, whereas the quenched statements are exact.

In pairwise voter dynamics, annealing the long-range links of a
small-world network preserves much of the
phenomenology~\cite{Vilone2004}. Here, by contrast, annealing at fixed
neighborhood size changes the stationary order parameter at order one in
$n$ within the computed range: the exact steady state is a property of
the quenched geometry, not of the neighborhood size alone.

% ============================================================
\section{Beyond one-step erasure}
\label{sec:regionII}
% ============================================================

\subsection{One-step extinction probability}

When R2 fails, a newly created minority is not necessarily erased by
the first fire of one of its nodes. For residuals $R_i$ define
\begin{align}
\cG&=\frac{1}{n^2}\sum_{i,j}\mathbf 1\{R_i\subseteq B_j\},\label{eq:cG}\\
\ctil&=\frac{\sum_i\sum_{j\in R_i}\mathbf 1\{R_i\subseteq B_j\}}{\sum_i|R_i|},
\label{eq:containment}
\end{align}
with $\ctil=1$ if all residuals are empty. $\cG$ is a purely geometric
overlap statistic. $\ctil$ is the probability that a newly created
minority (a residual) is erased by the first fire of one of its own
nodes, because that fire is uniform in the residual~\cite{SM}; under R2,
$\ctil=1$.

\begin{proposition}[Ring containment]
\label{prop:ring}
On the ring with $\beta=\min(n,4r+1)$,
$\cG=\min\{1,\max(0,2\beta-n+1)/n\}$ and
\begin{equation}
\ctil=
\begin{cases}
0, & 2\beta<n,\\
\dfrac{2\beta-n+1}{n-\beta}, & 2\beta\ge n,\ 3\beta<2n-1,\\[6pt]
1, & 3\beta\ge2n-1,
\end{cases}
\label{eq:ctil}
\end{equation}
continuous at $3\beta=2n-1$.
\end{proposition}

\begin{proof}
For $\beta<n$, a ball arc of length $\beta$ contains the residual arc of
length $n-\beta$ if and only if its center lies in an interval of
$2\beta-n+1$ positions; for $\ctil$ intersect this interval with the
residual. At $2\beta=n$ the residual is itself a ball and exactly one
node of it is admissible. Details in~\cite{SM}.
\end{proof}

Equation~\eqref{eq:ctil} agrees exactly, in rational arithmetic, with
exhaustive enumeration in all 192 cases $(n,r)$ with $12\le n\le40$
(Fig.~\ref{fig:containment}). At the Region~I/II boundary one-step extinction is lost: $\ctil$ falls
continuously below $1$ once the residual diameter exceeds the
interaction range, and it vanishes when R1 fails.

\begin{figure}[t]
\centering
\includegraphics[width=\columnwidth]{fig3_containment}
\caption{Ring containment for $n=40$. (a)~$\cG$ and (b)~$\ctil$ against
$q$: closed forms (Proposition~\ref{prop:ring}) at every integer
$\beta$ (steps) and exhaustive enumeration at the realized
$\beta=\min(n,4r+1)$ (circles). Dashed lines in (b): $\beta=n/2$ and
$\beta=(2n-1)/3$.}
\label{fig:containment}
\end{figure}

\subsection{Exact finite-size steady states in Region II}

For the ring in Region~II we solved the stationary two-state chain on
all configurations reachable from consensus. For the six odd $n\le25$
the reachable minority sets form a single irreducible class (73 to
51\,201 sets), the chain lumps exactly under the dihedral group and
state exchange, and the stationary values were obtained in exact
rational arithmetic and confirmed by an independent floating-point
solver of the unreduced chain (residual $\le4.1\times10^{-14}$). A third,
independent solver extends the computation to all 17 Region-II cases with
$n\le33$, odd and even (residual $<5\times10^{-13}$; Table~\ref{tab:exact}).

\begin{table*}[t]
\caption{All Region-II cases of the ring with $n\le33$
($n/2<\beta<(2n-1)/3$, $\beta=4r+1$). States: configurations reachable
from consensus, up to state exchange. $\E[M]$: exact stationary value of
the quenched dynamics; for the six odd $n\le25$ the exact rational is
given. $f(q)=(2-q)/(3-2q)$ is evaluated outside its domain as a
comparator. $\E[M]_{\rm ann}$: annealed version with the same $\beta$.}
\label{tab:exact}
\begin{ruledtabular}
\begin{tabular}{rrrcrrlccc}
$n$ & $r$ & $\beta$ & $q$ & $\ctil$ & states & $\E[M]$ & $f(q)$ & $\E[M]-f(q)$ & $\E[M]_{\rm ann}$\\
\colrule
 9 & 1 &  5 & 0.5556 & 1/2  & 73      & 0.754386 ($43/57$)   & 0.764706 & $-0.0103$ & 0.733333\\
15 & 2 &  9 & 0.6000 & 2/3  & 481     & 0.773810 ($65/84$)   & 0.777778 & $-0.0040$ & 0.750000\\
16 & 2 &  9 & 0.5625 & 3/7  & 1025    & 0.756831             & 0.766667 & $-0.0098$ & 0.732639\\
17 & 2 &  9 & 0.5294 & 1/4  & 2177    & 0.741294 ($149/201$) & 0.757576 & $-0.0163$ & 0.717647\\
21 & 3 & 13 & 0.6190 & 3/4  & 2689    & 0.781711 ($265/339$) & 0.783784 & $-0.0021$ & 0.757576\\
22 & 3 & 13 & 0.5909 & 5/9  & 5633    & 0.769553             & 0.775000 & $-0.0054$ & 0.744318\\
23 & 3 & 13 & 0.5652 & 2/5  & 11777   & 0.757895 ($72/95$)   & 0.767442 & $-0.0095$ & 0.732441\\
24 & 3 & 13 & 0.5417 & 3/11 & 24577   & 0.746914             & 0.760870 & $-0.0140$ & 0.721726\\
25 & 3 & 13 & 0.5200 & 1/6  & 51201   & 0.736721 ($319/433$) & 0.755102 & $-0.0184$ & 0.712000\\
27 & 4 & 17 & 0.6296 & 4/5  & 13825   & 0.785965             & 0.787234 & $-0.0013$ & 0.761905\\
28 & 4 & 17 & 0.6071 & 7/11 & 28673   & 0.776560             & 0.780000 & $-0.0034$ & 0.751190\\
29 & 4 & 17 & 0.5862 & 1/2  & 59393   & 0.767367             & 0.773585 & $-0.0062$ & 0.741379\\
30 & 4 & 17 & 0.5667 & 5/13 & 122881  & 0.758488             & 0.767857 & $-0.0094$ & 0.732353\\
31 & 4 & 17 & 0.5484 & 2/7  & 253953  & 0.750000             & 0.762712 & $-0.0127$ & 0.724014\\
32 & 4 & 17 & 0.5312 & 1/5  & 524289  & 0.741958             & 0.758065 & $-0.0161$ & 0.716283\\
33 & 4 & 17 & 0.5152 & 1/8  & 1081345 & 0.734396             & 0.753846 & $-0.0195$ & 0.709091\\
33 & 5 & 21 & 0.6364 & 5/6  & 67585   & 0.788618             & 0.789474 & $-0.0009$ & 0.764706\\
\end{tabular}
\end{ruledtabular}
\end{table*}

All 17 values lie below the closed form, by between $0.0009$ and
$0.0195$, and above the annealed values.

\begin{conjecture}
\label{conj:upper}
In Region~II of the ring, $\E[M]\le(2-q)/(3-2q)$ for all $n$ and all
admissible $r$.
\end{conjecture}

By Corollary~\ref{cor:duality}, which needs only R1, the conjecture is
equivalent to $\E[K]\ge(3-2q)/(2-q)$: below the R2 threshold, sign
reversals of $m$ would be rarer per update than in the exactly solvable
regime. Across the 17 cases the deviation increases monotonically with
$q$, from $-0.0195$ at $q=0.515$ to $-0.0009$ at $q=0.636$, whereas
$\ctil$ does not (it is $1/2$ at $q=0.556$ and $2/5$ at $q=0.565$); within
this set, $q$ organizes the deviation better than the one-step
quantity $\ctil$.

\emph{Interpretation.} In Region~I the largest minority that a single
update can create has diameter at most the interaction range, and the
closure of Sec.~\ref{sec:band} is exact. In Region~II a minority can
survive the first fire from inside it, possibly as several fragments,
and longer-lived minorities lower the time-averaged majority fraction.
The negative deviations are consistent with this picture, but the exact
chain gives the stationary mean, not a decomposition by mechanism, and
we do not isolate multi-step survival as the cause.

\subsection{Exploratory test across graph families}

At intermediate coverage, simulations of four graph families (rings,
two-dimensional lattices with eight neighbors, configuration-model and
Watts--Strogatz graphs) show family-dependent stationary values at equal
$(n,q)$. We tested whether $\ctil$ accounts for this dependence with a
weighted regression on 216 runs, using interaction sets truncated to a
common $\beta$. It does not: the coefficient of $\ctil$ is negative,
$-0.0196$ [cluster-bootstrap 95\% interval $(-0.0222,-0.0169)$], and
including $\ctil$ increases the separation between families. Because
$\ctil$ varies almost only within the ring family, the test does not
identify an alternative mechanism. The design, the analysis record and
all sensitivity analyses are given in the Supplemental
Material~\cite{SM}.

% ============================================================
\section{Discussion}
\label{sec:discussion}
% ============================================================

Group-flip dynamics on quenched neighborhoods provides an exactly
solvable nonequilibrium steady state without external noise. Under R1
and R2 with equal balls the stationary majority fraction, the mean
absolute magnetization, the rotation and erasure currents and the cycle
law are known in closed form for every finite $n$, and they depend on
the geometry only through the coverage $q$. The solution does not use a
moment closure: the conditional mean minority size is exactly affine,
because self-containment turns the overlap between a minority and a
random majority ball into a function of the minority size alone
[Eq.~\eqref{eq:overlap}]. The same mechanism explains why an annealed
description with the same group size fails: it discards the
correlation between minority domains and neighborhoods that the
quenched geometry creates.

\emph{Relation to annealed group dynamics.} Kim \emph{et
al.}~\cite{Kim2025} obtain leading-order prefactors of a logarithmic
exit time for a unanimity-based rule on annealed hypergraphs, with an
optimal nonlinearity that results from a competition between diffusive
and drift-dominated behavior. Their edge-update variant also changes several nodes at
once, but it copies a unanimous state and keeps consensus absorbing.
Here the moved set is a same-state part of a quenched ball, consensus is
not absorbing, and the object computed is a stationary value rather than
an exit time. The two results therefore address different questions; the
annealed comparison of Sec.~\ref{sec:annealed} shows that, for the
group-flip rule, replacing the quenched balls by annealed sets of the
same size changes the stationary state.

\emph{Limitations.} (i)~All exact results are for two states.
(ii)~The closed form~\eqref{eq:band} requires equal balls; for unequal
balls Proposition~\ref{prop:het} gives an exact expression with one
effective parameter and exact bounds. (iii)~R1 forces the radius to grow
as $n^{1/d}$, so the large-$n$ limit at fixed $q$ is not a short-range
limit; Theorem~\ref{thm:band} is a finite-$n$ statement about extensive,
spatially fixed neighborhoods. (iv)~Region~II is solved only
numerically, for $n\le33$, and Conjecture~\ref{conj:upper} is unproven.
(v)~The annealed size dependence is computed, not proved.

\emph{Open problems.} A theory of Region~II must describe minorities that
survive their first erasure attempt, for instance through multi-step
containment probabilities or the statistics of minority fragments. The
fluctuations of $M$ are not fixed by the present closure, since they
involve pairwise overlaps of balls. A proof that $q^\ast\le\bar q$ for
unequal balls would turn the observation of Sec.~\ref{sec:band} into a
bound. Finally, whether a short-range version of group-flip dynamics
retains an ordered, sign-reversing steady state is open.

\begin{acknowledgments}
During the preparation of this work the author used Claude Opus 5.5 and
Claude Sonnet 5.5 (Anthropic) for code development, LaTeX editing,
manuscript review and language editing, and Qwen 3.8 Max (Alibaba) for
literature search. The author reviewed and edited all content and takes
full responsibility for it. No funding was received for this study.
\end{acknowledgments}

\section*{Data availability}
The code and data that support the findings of this article are openly
available~\cite{repo}, at
\url{https://doi.org/10.5281/zenodo.24050030} and
\url{https://github.com/phmldtngyn/Group-flip-consensus-on-hypergraphs}.

\appendix

% ============================================================
\section{Ergodic step in Theorem~\ref{thm:band}}
\label{app:ergodic}
% ============================================================

Start the epoch chain in an invariant distribution $\pi$, which exists
because the state space is finite. Let $Y_m=\int_{\text{cycle }m}M\,dt$
and $\tau_m$ be the reward and length of cycle $m$. The sequence
$(Y_m,\tau_m)$ is stationary as a functional of a stationary chain; it is
not independent, and independence is not used. By the Birkhoff ergodic
theorem, $N^{-1}\sum_{m\le N}Y_m\to\E[Y\mid\mathcal I]$ and
$N^{-1}\sum_{m\le N}\tau_m\to\E[\tau\mid\mathcal I]\ge1/n$ almost
surely, with $\mathcal I$ the invariant $\sigma$-field. On the event that
the chain lies in a closed class $c$, these are the class means
$\E_c[Y]=1/n$ and $\E_c[\tau]=(1+\bar x_c)/n$. Applying
Eq.~\eqref{eq:xbar} to the chain restricted to $c$ gives
$\bar x_c=(1-q)/(2-q)$ for every class, so the ratio is $1/(1+\bar x)$
almost surely. From an arbitrary initial condition the epoch chain
enters a closed class in finite time almost surely, so the long-run
time average of $M$ is the same for every initial law~\cite{SM}.

\bibliography{refs_pre}

\end{document}
